\documentclass[11pt,a4paper]{article}

\usepackage[utf8]{inputenc}
\usepackage[T1]{fontenc}
\usepackage{amsmath,amssymb,amsthm}
\usepackage{geometry}
\usepackage{booktabs}
\usepackage{hyperref}
\usepackage{cleveref}

\geometry{margin=2.5cm}

\newcommand{\dd}{\mathrm{d}}
\newcommand{\pp}{p^2}
\newcommand{\kk}{k^2}
\newcommand{\TT}{\mathrm{TT}}
\newcommand{\tr}{\mathrm{tr}}
\newcommand{\Lap}{\Delta}

\newtheorem{theorem}{Theorem}
\newtheorem{corollary}{Corollary}
\newtheorem{proposition}{Proposition}
\theoremstyle{definition}
\newtheorem{definition}{Definition}
\theoremstyle{remark}
\newtheorem{remark}{Remark}

\title{Position-Space Correlation Functions\\
for Classical Stochastic Gravity\\[4pt]
$\boldsymbol{S = \alpha\, R^2 - \beta\, R_{\mu\nu}\,R^{\mu\nu}}$}
\author{Computed with TensorGR.jl}
\date{\today}

\begin{document}
\maketitle

\begin{abstract}
We compute the complete set of two-point correlation functions for
the classical metric perturbation $h_{\mu\nu}$ governed by the
Onsager--Machlup action
$S = \alpha\,R^2 - \beta\,R_{\mu\nu}R^{\mu\nu}$
linearised around flat space.
All results are expressed in position space as Green's functions
of fourth-order differential operators.
The metric $h_{\mu\nu}$ is a real classical stochastic field:
no ghosts, no unitarity constraints, no $i\varepsilon$~prescriptions
arise.  The path integral $\int\!\mathcal{D}h\;e^{-S}$ defines a
probability measure on metric configurations whose covariance is
the inverse kinetic operator.
The momentum-space kinetic matrices and their inverses are verified
numerically with \texttt{TensorGR.jl}; the position-space Green's
functions are standard analytic results.
\end{abstract}

\tableofcontents

%% ======================================================================
\section{Classical stochastic gravity}
\label{sec:setup}

\subsection{The Onsager--Machlup framework}

A classical metric $g_{\mu\nu}$ subject to stochastic fluctuations is
described by the probability functional
\begin{equation}\label{eq:OM}
  P[g] \;\propto\; e^{-S[g]}\,,
  \qquad
  S[g] = \int \dd^4 x\;\sqrt{|g|}\;\bigl[\,
    \alpha\,R^2 - \beta\,R_{\mu\nu}\,R^{\mu\nu}
  \,\bigr]\,.
\end{equation}
This is an Onsager--Machlup action: $S$ is real and bounded below, so
that~\eqref{eq:OM} is a normalisable measure on the space of metrics.
The coefficients $\alpha,\beta$ encode the noise strength of whatever
stochastic dynamics (Langevin equation, fluctuation--dissipation relation)
drives the metric fluctuations.

The two-point correlation function of the classical metric perturbation
$h_{\mu\nu} = g_{\mu\nu} - \eta_{\mu\nu}$ is
\begin{equation}\label{eq:corr-def}
  C_{\mu\nu\rho\sigma}(x,y)
  \equiv \langle h_{\mu\nu}(x)\, h_{\rho\sigma}(y)\rangle
  = \frac{1}{Z}\int\!\mathcal{D}h\; h_{\mu\nu}(x)\,h_{\rho\sigma}(y)\;
    e^{-S^{(2)}[h]}\,,
\end{equation}
where $S^{(2)}$ is the action expanded to quadratic order.
For a Gaussian theory this gives
\begin{equation}\label{eq:gauss}
  C = \tfrac{1}{2}\,\mathcal{M}^{-1}\,,
\end{equation}
with $\mathcal{M}$ the kinetic operator in
$S^{(2)} = \int h\,\mathcal{M}\, h\;\dd^4x$.
The factor $\frac{1}{2}$ is the standard Gaussian normalisation;
it can be absorbed into a redefinition of $\alpha,\beta$.
In what follows we compute $G = \mathcal{M}^{-1}$ (the Green's function
of the kinetic operator) and note that
$C = \tfrac{1}{2}\,G$.

\paragraph{No ghosts, no $i\varepsilon$.}
Unlike in quantum gravity, the metric here is a \emph{real classical
field}.  The higher-derivative kinetic operator $\Box^2$ does not
produce ghost degrees of freedom, because there is no Hilbert space
and no unitarity requirement.  The $1/p^4$ correlation functions
are simply the covariance of a smooth stochastic field with
fourth-order noise.

\subsection{Conventions}

We expand $g_{\mu\nu} = \eta_{\mu\nu} + h_{\mu\nu}$ with signature
$(-,+,+,+)$.  The d'Alembertian is
$\Box = -\partial_t^2 + \nabla^2$, and the spatial Laplacian is
$\nabla^2 = \partial_i\partial_i$.

\paragraph{Bardeen gauge.}
The SVT decomposition in Bardeen gauge gives
\begin{equation}\label{eq:svt}
  h_{00} = 2\Phi\,,\qquad
  h_{0i} = V_i\,,\qquad
  h_{ij} = 2\psi\,\delta_{ij} + h^\TT_{ij}\,,
\end{equation}
with $\partial^i V_i = 0$ (transverse vector) and
$\partial^i h^\TT_{ij} = 0$, $h^\TT_{ii} = 0$
(transverse-traceless tensor).
The trace is $h = -2\Phi + 6\psi$.

%% ======================================================================
\section{Position-space action}
\label{sec:action}

\subsection{Linearised curvature as differential operators}

The linearised Ricci tensor on a flat background is a second-order
differential operator acting on $h_{\mu\nu}$:
\begin{equation}\label{eq:lin-ricci}
  R^{(1)}_{\mu\nu}[h]
  = \tfrac{1}{2}\bigl(
      \partial^\rho\partial_\mu h_{\rho\nu}
    + \partial^\rho\partial_\nu h_{\rho\mu}
    - \Box\, h_{\mu\nu}
    - \partial_\mu\partial_\nu h
  \bigr)\,.
\end{equation}
The linearised Ricci scalar is
\begin{equation}\label{eq:lin-R}
  R^{(1)}[h] = \partial_\mu\partial_\nu h^{\mu\nu} - \Box\, h\,.
\end{equation}
These are \emph{local} differential expressions in position space.
The quadratic action is therefore
\begin{equation}\label{eq:action-pos}
  S^{(2)}[h] = \int\dd^4x\;\Bigl[\,
    \alpha\bigl(\partial_\mu\partial_\nu h^{\mu\nu} - \Box h\bigr)^2
  - \frac{\beta}{4}
    \bigl(\partial^\rho\partial_\mu h_{\rho\nu}
         + \partial^\rho\partial_\nu h_{\rho\mu}
         - \Box h_{\mu\nu}
         - \partial_\mu\partial_\nu h\bigr)^2
  \,\Bigr]\,,
\end{equation}
a manifestly local functional involving up to four derivatives of
$h_{\mu\nu}$.

\subsection{Ricci components on SVT fields}

Evaluating~\eqref{eq:lin-ricci} on the Bardeen decomposition~\eqref{eq:svt}
gives (see \cref{tab:ricci}):

\begin{table}[ht]
\centering
\caption{Linearised Ricci components on SVT fields.  All expressions
are position-space differential operators.}
\label{tab:ricci}
\begin{tabular}{@{}ll@{}}
\toprule
Component & Position-space expression \\
\midrule
$R_{00}$ & $-\nabla^2\Phi - 3\,\ddot\psi$ \\[3pt]
$R_{0i}\big|_V$ & $-\frac{1}{2}\nabla^2 V_i$ \\[3pt]
$R_{0i}\big|_S$ & $-2\,\partial_i\dot\psi$ \\[3pt]
$R_{ij}\big|_S$ & $\partial_i\partial_j(\Phi-\psi) - \Box\psi\;\delta_{ij}$ \\[3pt]
$R_{ij}\big|_\TT$ & $-\frac{1}{2}\Box\, h^\TT_{ij}$ \\[3pt]
$R^{(1)}$ & $2\nabla^2\Phi - 4\nabla^2\psi + 6\ddot\psi$ \\
\bottomrule
\end{tabular}
\end{table}

Here $\dot f = \partial_t f$ and
$\ddot f = \partial_t^2 f$.  These are the same as the Fourier-space
expressions under the correspondence
$\omega^2 \leftrightarrow -\partial_t^2$,
$k^2 \leftrightarrow -\nabla^2$,
$p^2 \leftrightarrow \Box$.

%% ======================================================================
\section{Kinetic operators by sector}
\label{sec:operators}

Substituting the Ricci components into $S^{(2)} = \alpha(R^{(1)})^2
- \beta\,R^{(1)}_{\mu\nu}R^{(1)\mu\nu}$ yields a fourth-order
differential operator for each SVT sector.

\subsection{Tensor sector}

Since $R^{(1)}|_\TT = 0$ and
$R^{(1)}_{ij}|_\TT = -\frac{1}{2}\Box h^\TT_{ij}$, the tensor
action is
\begin{equation}\label{eq:S-TT}
  \boxed{
  S_\TT = -\frac{\beta}{4}\int\dd^4x\;
    h^\TT_{ij}\;\Box^2\; h^{\TT\,ij}\,.}
\end{equation}
The kinetic operator is $\mathcal{M}_\TT = -\frac{\beta}{4}\,\Box^2$.

\subsection{Vector sector}

Since $R^{(1)}|_V = 0$ and
$R^{(1)}_{0i}|_V = -\frac{1}{2}\nabla^2 V_i$, the vector action is
\begin{equation}\label{eq:S-V}
  \boxed{
  S_V = \frac{\beta}{2}\int\dd^4x\;
    V_i\;\nabla^2\Box\; V^i\,.}
\end{equation}
The kinetic operator is
$\mathcal{M}_V = \frac{\beta}{2}\,\nabla^2\Box = -\frac{\beta}{2}(-\nabla^2)(-\Box)$.

\subsection{Scalar sector}

The scalar action couples $\Phi$ and $\psi$ through a $2\times 2$
matrix of differential operators.  Define
\begin{equation}\label{eq:AB}
  A \equiv 2\alpha - \beta\,,\qquad
  B \equiv 3\alpha - \beta\,.
\end{equation}

\begin{theorem}[Scalar kinetic operator]\label{thm:scalar-M}
The scalar sector of the quadratic action is
$S_S = \int\dd^4x\;\phi^I\,\mathcal{M}_{IJ}\,\phi^J$
with $\phi = (\Phi,\psi)^T$ and
\begin{equation}\label{eq:scalar-M}
  \mathcal{M} = \begin{pmatrix}
    2A\,\nabla^4
    & -4B\,\nabla^2\Box + 2A\,\nabla^4 \\[6pt]
    -4B\,\nabla^2\Box + 2A\,\nabla^4
    & 12B\,\Box^2 - 8B\,\nabla^2\Box + 2A\,\nabla^4
  \end{pmatrix},
\end{equation}
where $\nabla^4 \equiv (\nabla^2)^2$.
\end{theorem}

The correspondence with the Fourier-space kinetic matrix is
$\nabla^4 \leftrightarrow k^4$,
$-\nabla^2\Box \leftrightarrow k^2 p^2$,
$\Box^2 \leftrightarrow p^4$.

\begin{remark}
The three differential operators $\nabla^4$, $\nabla^2\Box$, and
$\Box^2$ are all fourth-order, but have different symmetry:
$\Box^2$ is Lorentz-invariant,
$\nabla^4$ is purely spatial (instantaneous), and
$\nabla^2\Box$ is mixed.
This trichotomy governs the structure of the correlation functions.
\end{remark}

%% ======================================================================
\section{Correlation functions}
\label{sec:greens}

The two-point correlation function of the classical metric is
$C = \frac{1}{2}\mathcal{M}^{-1}$,
i.e.\ half the Green's function of the kinetic operator
(\cref{eq:gauss}).  We compute $G = \mathcal{M}^{-1}$.

\subsection{Building-block Green's functions}
\label{sec:building}

All correlation functions decompose into three building blocks, the
Green's functions of the three fourth-order operators
$\Box^2$, $(-\nabla^2)(-\Box)$, and $\nabla^4$.

\begin{definition}[Elementary Green's functions]\label{def:greens}
\begin{alignat}{2}
  &\Box^2\; G_{\Box^2}(x)
  &&= \delta^4(x)\,,
  \label{eq:G-box2}\\[4pt]
  &(-\nabla^2)(-\Box)\; G_{\nabla^2\Box}(x)
  &&= \delta^4(x)\,,
  \label{eq:G-mix}\\[4pt]
  &\nabla^4\; G_{\nabla^4}(x)
  &&= \delta^4(x)\,.
  \label{eq:G-lap4}
\end{alignat}
\end{definition}

The biharmonic Green's function $G_{\Box^2}$ is the most regular;
it exists as a locally integrable function.  The spatial biharmonic
$G_{\nabla^4}$ is singular: it has no time dependence (the operator
$\nabla^4$ does not involve $\partial_t$) and produces an
\emph{instantaneous} correlation.

For the Euclidean theory (see \cref{sec:euclidean}), these have
explicit closed forms.

\subsection{Tensor correlation function}

The Green's function equation for the tensor sector is
\begin{equation}
  -\frac{\beta}{4}\;\Box^2\;
  G^\TT_{ijkl}(x-y) = \Pi^\TT_{ijkl}\;\delta^4(x-y)\,,
\end{equation}
where $\Pi^\TT_{ijkl} = \frac{1}{2}(P^\mathrm{T}_{ik}P^\mathrm{T}_{jl}
+ P^\mathrm{T}_{il}P^\mathrm{T}_{jk} - P^\mathrm{T}_{ij}P^\mathrm{T}_{kl})$
with $P^\mathrm{T}_{ij} = \delta_{ij} - \partial_i\partial_j/\nabla^2$.
The solution is
\begin{equation}\label{eq:corr-TT}
  \boxed{
  \bigl\langle h^\TT_{ij}(x)\, h^\TT_{kl}(y)\bigr\rangle
  = -\frac{2}{\beta}\;
    \Pi^\TT_{ijkl}\; G_{\Box^2}(x-y)\,.}
\end{equation}

\subsection{Vector correlation function}

The Green's function equation is
\begin{equation}
  \frac{\beta}{2}\;\nabla^2\Box\;
  G^V_{ij}(x-y) = P^\mathrm{T}_{ij}\;\delta^4(x-y)\,.
\end{equation}
The solution is
\begin{equation}\label{eq:corr-V}
  \boxed{
  \bigl\langle V_i(x)\, V_j(y)\bigr\rangle
  = -\frac{1}{\beta}\;
    P^\mathrm{T}_{ij}\; G_{\nabla^2\Box}(x-y)\,.}
\end{equation}

\subsection{Scalar correlation function}

The scalar Green's function $G^S_{IJ}(x-y)$ satisfies
$\mathcal{M}_{IK}\,G^S_{KJ}(x-y) = \delta_{IJ}\,\delta^4(x-y)$.
Its entries decompose as linear combinations of the three building
blocks.

\begin{theorem}[Scalar correlators]\label{thm:scalar-corr}
\begin{equation}\label{eq:corr-scalar}
  \boxed{
  \begin{aligned}
    \langle\Phi(x)\,\Phi(y)\rangle
    &= -\frac{3}{2\beta}\;G_{\nabla^4}(x-y)
     - \frac{1}{\beta}\;G_{\nabla^2\Box}(x-y)
     - \frac{A}{4\beta B}\;G_{\Box^2}(x-y)\,,\\[6pt]
    \langle\Phi(x)\,\psi(y)\rangle
    &= \frac{1}{2\beta}\;G_{\nabla^2\Box}(x-y)
     + \frac{A}{4\beta B}\;G_{\Box^2}(x-y)\,,\\[6pt]
    \langle\psi(x)\,\psi(y)\rangle
    &= -\frac{A}{4\beta B}\;G_{\Box^2}(x-y)\,,
  \end{aligned}}
\end{equation}
where $A = 2\alpha - \beta$ and $B = 3\alpha - \beta$.
\end{theorem}

\begin{proof}
Invert the $2\times2$ kinetic matrix in Fourier space
(Cramer's rule with $\det\mathcal{M} = -8\beta B\,k^4 p^4$),
then identify each monomial $k^{-4}$, $(k^2 p^2)^{-1}$, $p^{-4}$
with the corresponding Green's function.
\end{proof}

\begin{remark}[Structure of the scalar correlator]
The three terms have distinct physical character:
\begin{enumerate}
\item The $G_{\nabla^4}$ term is \emph{instantaneous}: it
  represents a spatial correlation that is independent of the time
  separation $|t_x - t_y|$.  It appears only in
  $\langle\Phi\Phi\rangle$ and reflects the constraint character
  of the lapse perturbation.
\item The $G_{\nabla^2\Box}$ term mixes spatial and temporal
  propagation.  It appears in $\langle\Phi\Phi\rangle$ and
  $\langle\Phi\psi\rangle$ but not in $\langle\psi\psi\rangle$.
\item The $G_{\Box^2}$ term is fully covariant (Lorentz-invariant
  operator) and appears in all three entries.
  It is the only contribution to $\langle\psi\psi\rangle$ and
  encodes the genuinely propagating scalar degree of freedom.
\end{enumerate}
\end{remark}

%% ======================================================================
\section{Euclidean Green's functions}
\label{sec:euclidean}

For a classical stochastic field the path integral $\int\!\mathcal{D}h\;
e^{-S}$ must converge, which requires positive-definite kinetic operators.
This is naturally achieved in the Euclidean formulation obtained by
Wick-rotating $t \to -i\tau$.  The Lorentzian operators become
Euclidean elliptic operators:
\begin{equation}
  \Box \;\to\; -\Lap_E\,,\qquad
  \Lap_E = \partial_\tau^2 + \nabla^2\,,\qquad
  |x_E|^2 = \tau^2 + |\mathbf{x}|^2\,.
\end{equation}
The Euclidean kinetic operators are positive-definite for
$\alpha > 0$ and $\beta > 0$.

\subsection{Biharmonic Green's function}

The Green's function of $\Lap_E^2$ in four Euclidean dimensions is
logarithmic (critical dimension):
\begin{equation}\label{eq:GE-biharm}
  \boxed{
  \Lap_E^2\; G_{\Box^2}^E(x_E) = \delta^4(x_E)
  \qquad\Longrightarrow\qquad
  G_{\Box^2}^E(x_E) = -\frac{1}{8\pi^2}\,
    \log\frac{|x_E|^2}{\ell^2}\,,}
\end{equation}
where $\ell$ is an infrared reference scale (the 4D biharmonic has a
logarithmic zero mode).  This is $O(4)$-invariant and smooth away
from the origin.

\subsection{Spatial biharmonic Green's function}

The operator $\nabla^4$ acts only on the three spatial coordinates,
so its Green's function is instantaneous in Euclidean time:
\begin{equation}\label{eq:GE-splap4}
  \boxed{
  \nabla^4\; G_{\nabla^4}^E(\tau, \mathbf{x})
  = \delta(\tau)\,\delta^3(\mathbf{x})
  \qquad\Longrightarrow\qquad
  G_{\nabla^4}^E(\tau,\mathbf{x})
  = -\frac{|\mathbf{x}|}{8\pi}\;\delta(\tau)\,.}
\end{equation}
This is the 3D biharmonic Green's function
($\nabla^4 G = \delta^3$ in $\mathbb{R}^3$ gives
$G(r) = -r/(8\pi)$) times a delta function in Euclidean time.

\subsection{Mixed Green's function}

The mixed operator $(-\nabla^2)(-\Lap_E)$ factorises into a
spatial and a full Laplacian, but the factors share the spatial
coordinates.  The Green's function is obtained by integrating out
the temporal Fourier mode first:
\begin{equation}\label{eq:GE-mix}
  \boxed{
  G_{\nabla^2\Box}^E(\tau, \mathbf{x})
  = \frac{1}{4\pi^2}\int_0^\infty
    \frac{e^{-\kappa|\tau|}\,\sin(\kappa\,|\mathbf{x}|)}
         {\kappa^2\,|\mathbf{x}|}\;
    \frac{\dd\kappa}{\pi}\,,}
\end{equation}
where the $\kappa$~integral runs over the magnitude of the spatial
momentum.  At equal times ($\tau = 0$):
\begin{equation}\label{eq:GE-mix-et}
  G_{\nabla^2\Box}^E(0, \mathbf{x})
  = \frac{1}{4\pi^2}\int_0^\infty
    \frac{\sin(\kappa\,|\mathbf{x}|)}
         {\kappa^2\,|\mathbf{x}|}\;
    \frac{\dd\kappa}{\pi}
  = \frac{1}{4\pi^2\,|\mathbf{x}|}\cdot\frac{1}{2}
  = \frac{1}{8\pi^2\,|\mathbf{x}|}\,.
\end{equation}

For large Euclidean separation $|x_E| \gg \ell$, the mixed Green's
function approaches the biharmonic one up to anisotropic corrections:
$G_{\nabla^2\Box}^E(x_E) \sim -\frac{1}{8\pi^2}\log(|x_E|/\ell)
+ \mathcal{O}(1)$.

%% ======================================================================
\section{Equal-time spatial correlators}
\label{sec:equal-time}

The most directly observable quantities for a classical stochastic metric
are the equal-time spatial correlation functions
$C(\mathbf{r}) = \langle h(t,\mathbf{x})\, h(t,\mathbf{x}+\mathbf{r})\rangle$.
These are obtained by evaluating the Green's functions at $\tau = 0$.

\subsection{Equal-time power spectra}

In Fourier space (spatial momentum $\mathbf{k}$, integrating over
frequency), the equal-time power spectra for each sector are:

\begin{proposition}[Equal-time power spectra]\label{prop:power}
In the Euclidean formulation, integrating
$\int\!\frac{\dd\omega_E}{2\pi}$ over the Euclidean frequency gives
the equal-time spectral densities
\begin{align}
  P_\TT(k) &= \frac{1}{\beta\, k^3}\,,
    \label{eq:P-TT}\\[4pt]
  P_V(k) &= \frac{1}{\beta\, k^3}\,,
    \label{eq:P-V}\\[4pt]
  P_{\psi\psi}(k) &= \frac{A}{4\beta B}\cdot\frac{1}{k^3}\,,
    \label{eq:P-psi}
\end{align}
per polarisation, where $k = |\mathbf{k}|$.
\end{proposition}

\begin{proof}
For the tensor sector, the Euclidean propagator is
$(4/\beta)\,(\omega_E^2 + k^2)^{-2}$.  Integrating:
$\int \frac{\dd\omega_E}{2\pi}\,\frac{4}{\beta(\omega_E^2+k^2)^2}
= \frac{4}{\beta}\cdot\frac{1}{4k^3} = \frac{1}{\beta k^3}$.
The vector and scalar follow analogously, using
$\int\frac{\dd\omega_E}{2\pi}\,(\omega_E^2+k^2)^{-1} = \frac{1}{2k}$
for the $1/(k^2 p_E^2)$ poles.
\end{proof}

\begin{remark}[Scale-invariant spectrum]
All sectors share the same spectral shape $P(k) \propto k^{-3}$.
The equal-time variance diverges logarithmically:
\begin{equation}
  \langle h^2 \rangle
  \propto \int \frac{\dd^3 k}{(2\pi)^3}\;P(k)
  = \frac{1}{2\pi^2}\int_0^\infty \frac{\dd k}{k}
  = \frac{1}{2\pi^2}\,\log\frac{k_{\max}}{k_{\min}}\,.
\end{equation}
This is the hallmark of a marginally scale-invariant stochastic field
in three spatial dimensions, analogous to a Harrison--Zel'dovich
spectrum.  The logarithmic divergence is cut off by the physical
UV and IR scales of the problem.
\end{remark}

\subsection{Equal-time spatial correlators}

From the power spectra, the 3D Fourier transforms give the
equal-time spatial correlation functions.  For the propagating
(non-instantaneous) contributions:

\begin{proposition}[Equal-time correlators]\label{prop:et}
The equal-time spatial correlators arising from the frequency-integrated
$1/p_E^4$ and $1/(k^2 p_E^2)$ poles are both logarithmic in the
spatial separation $r = |\mathbf{r}|$:
\begin{equation}\label{eq:et-log}
  \int\frac{\dd^3k}{(2\pi)^3}\;\frac{e^{i\mathbf{k}\cdot\mathbf{r}}}{k^3}
  = -\frac{1}{2\pi^2}\,\log(\mu\, r) + \text{const}\,,
\end{equation}
where $\mu$ is an infrared cutoff.
\end{proposition}

Combining with the instantaneous $G_{\nabla^4}$ contribution from
\cref{eq:GE-splap4}, the full equal-time scalar correlator is:
\begin{align}
  \langle\Phi(\mathbf{x})\,\Phi(\mathbf{x}+\mathbf{r})\rangle
  \big|_{t=t'}
  &= \frac{3}{16\pi\beta}\;r\;\delta(0)
   + \frac{1}{2\pi^2\beta}\;\log(\mu r)
   + \frac{A}{8\pi^2\beta B}\;\log(\mu r)
   + \text{const}\,,
  \label{eq:et-PhiPhi}
\end{align}
where the $\delta(0)$ in the first term reflects the white-noise
character of the instantaneous contribution and must be regularised
by the physical time resolution of the problem.

%% ======================================================================
\section{Special limits}
\label{sec:limits}

\subsection{Conformal point: $\beta = 3\alpha$ \texorpdfstring{$(B=0)$}{(B=0)}}

The action becomes $S = -3\alpha\,C_{\mu\nu\rho\sigma}C^{\mu\nu\rho\sigma}$
(linearised Weyl-squared).  Conformal symmetry eliminates the scalar
degree of freedom: $\det\mathcal{M}_S = 0$, and $\langle\psi\psi\rangle$
diverges.  Only the tensor and vector correlators survive, with the
\emph{same} Green's function structure as the general case.

\subsection{Trace-decoupling: $\beta = 2\alpha$ \texorpdfstring{$(A=0)$}{(A=0)}}

The $G_{\Box^2}$ contribution drops out of all scalar correlators.
In particular $\langle\psi\psi\rangle = 0$: the conformal factor
has no fluctuations.  Only the mixed and instantaneous terms remain
in $\langle\Phi\Phi\rangle$ and $\langle\Phi\psi\rangle$.

\subsection{Pure $R^2$: $\beta = 0$}

The tensor and vector kinetic operators vanish:
$\mathcal{M}_\TT = \mathcal{M}_V = 0$.
These modes are not damped by the noise and do not acquire
well-defined correlation functions.  The scalar determinant also
vanishes, reflecting the enhanced (conformal) gauge symmetry.

\subsection{Pure $R_{\mu\nu}R^{\mu\nu}$: $\alpha = 0$}

With $A = B = -\beta$:
\begin{equation}
  \langle\psi\psi\rangle\big|_{\alpha=0}
  = -\frac{1}{4\beta^2}\;G_{\Box^2}(x-y)\,.
\end{equation}
All correlators scale as $1/\beta$ or $1/\beta^2$.  The scalar sector
has the simplest structure, with $\Phi$ and $\psi$ entering
symmetrically.

%% ======================================================================
\section{Summary}
\label{sec:summary}

For the Onsager--Machlup action
$S = \alpha\,R^2 - \beta\,R_{\mu\nu}R^{\mu\nu}$
of a classical stochastic metric linearised around flat space,
the correlation functions are:

\medskip
\noindent\textbf{Kinetic operators} (position space):
\begin{equation*}
  \mathcal{M}_\TT = -\frac{\beta}{4}\,\Box^2\,,\qquad
  \mathcal{M}_V = \frac{\beta}{2}\,\nabla^2\Box\,,\qquad
  \mathcal{M}_S = \text{eq.~\eqref{eq:scalar-M}}\,.
\end{equation*}

\noindent\textbf{Tensor} (2~d.o.f.):
\begin{equation*}
  \bigl\langle h^\TT_{ij}(x)\, h^\TT_{kl}(y)\bigr\rangle
  = -\frac{2}{\beta}\;\Pi^\TT_{ijkl}\;G_{\Box^2}(x-y)
\end{equation*}

\noindent\textbf{Vector} (2~d.o.f.):
\begin{equation*}
  \bigl\langle V_i(x)\, V_j(y)\bigr\rangle
  = -\frac{1}{\beta}\;P^\mathrm{T}_{ij}\;G_{\nabla^2\Box}(x-y)
\end{equation*}

\noindent\textbf{Scalar} (2~coupled d.o.f., $A = 2\alpha-\beta$,
$B = 3\alpha-\beta$):
\begin{equation*}
  \begin{aligned}
  \langle\Phi\Phi\rangle
  &= -\tfrac{3}{2\beta}\,G_{\nabla^4}
   - \tfrac{1}{\beta}\,G_{\nabla^2\Box}
   - \tfrac{A}{4\beta B}\,G_{\Box^2}\,,\\[4pt]
  \langle\Phi\psi\rangle
  &= \tfrac{1}{2\beta}\,G_{\nabla^2\Box}
   + \tfrac{A}{4\beta B}\,G_{\Box^2}\,,\\[4pt]
  \langle\psi\psi\rangle
  &= -\tfrac{A}{4\beta B}\,G_{\Box^2}\,,
  \end{aligned}
\end{equation*}
where all Green's functions are evaluated at $x-y$.

\medskip
\noindent\textbf{Euclidean Green's functions} (4D):
\begin{equation*}
  G_{\Box^2}^E = -\frac{\log(|x_E|^2/\ell^2)}{8\pi^2}\,,\qquad
  G_{\nabla^4}^E = -\frac{|\mathbf{x}|}{8\pi}\,\delta(\tau)\,,\qquad
  G_{\nabla^2\Box}^E(0,\mathbf{x}) = \frac{1}{8\pi^2\,|\mathbf{x}|}\,.
\end{equation*}

\noindent\textbf{Equal-time power spectrum} (all sectors):
$P(k) \propto k^{-3}$ --- a marginally scale-invariant, logarithmically
correlated classical metric.

\end{document}
